% Options for packages loaded elsewhere
\PassOptionsToPackage{unicode}{hyperref}
\PassOptionsToPackage{hyphens}{url}
\PassOptionsToPackage{dvipsnames,svgnames,x11names}{xcolor}
%
\documentclass[
  11pt,
]{article}
\usepackage{amsmath,amssymb}
\usepackage{iftex}
\ifPDFTeX
  \usepackage[T1]{fontenc}
  \usepackage[utf8]{inputenc}
  \usepackage{textcomp} % provide euro and other symbols
\else % if luatex or xetex
  \usepackage{unicode-math} % this also loads fontspec
  \defaultfontfeatures{Scale=MatchLowercase}
  \defaultfontfeatures[\rmfamily]{Ligatures=TeX,Scale=1}
\fi
\usepackage{lmodern}
\ifPDFTeX\else
  % xetex/luatex font selection
\fi
% Use upquote if available, for straight quotes in verbatim environments
\IfFileExists{upquote.sty}{\usepackage{upquote}}{}
\IfFileExists{microtype.sty}{% use microtype if available
  \usepackage[]{microtype}
  \UseMicrotypeSet[protrusion]{basicmath} % disable protrusion for tt fonts
}{}
\makeatletter
\@ifundefined{KOMAClassName}{% if non-KOMA class
  \IfFileExists{parskip.sty}{%
    \usepackage{parskip}
  }{% else
    \setlength{\parindent}{0pt}
    \setlength{\parskip}{6pt plus 2pt minus 1pt}}
}{% if KOMA class
  \KOMAoptions{parskip=half}}
\makeatother
\usepackage{xcolor}
\usepackage[margin=1in]{geometry}
\usepackage{color}
\usepackage{fancyvrb}
\newcommand{\VerbBar}{|}
\newcommand{\VERB}{\Verb[commandchars=\\\{\}]}
\DefineVerbatimEnvironment{Highlighting}{Verbatim}{commandchars=\\\{\}}
% Add ',fontsize=\small' for more characters per line
\newenvironment{Shaded}{}{}
\newcommand{\AlertTok}[1]{\textcolor[rgb]{1.00,0.00,0.00}{\textbf{#1}}}
\newcommand{\AnnotationTok}[1]{\textcolor[rgb]{0.38,0.63,0.69}{\textbf{\textit{#1}}}}
\newcommand{\AttributeTok}[1]{\textcolor[rgb]{0.49,0.56,0.16}{#1}}
\newcommand{\BaseNTok}[1]{\textcolor[rgb]{0.25,0.63,0.44}{#1}}
\newcommand{\BuiltInTok}[1]{\textcolor[rgb]{0.00,0.50,0.00}{#1}}
\newcommand{\CharTok}[1]{\textcolor[rgb]{0.25,0.44,0.63}{#1}}
\newcommand{\CommentTok}[1]{\textcolor[rgb]{0.38,0.63,0.69}{\textit{#1}}}
\newcommand{\CommentVarTok}[1]{\textcolor[rgb]{0.38,0.63,0.69}{\textbf{\textit{#1}}}}
\newcommand{\ConstantTok}[1]{\textcolor[rgb]{0.53,0.00,0.00}{#1}}
\newcommand{\ControlFlowTok}[1]{\textcolor[rgb]{0.00,0.44,0.13}{\textbf{#1}}}
\newcommand{\DataTypeTok}[1]{\textcolor[rgb]{0.56,0.13,0.00}{#1}}
\newcommand{\DecValTok}[1]{\textcolor[rgb]{0.25,0.63,0.44}{#1}}
\newcommand{\DocumentationTok}[1]{\textcolor[rgb]{0.73,0.13,0.13}{\textit{#1}}}
\newcommand{\ErrorTok}[1]{\textcolor[rgb]{1.00,0.00,0.00}{\textbf{#1}}}
\newcommand{\ExtensionTok}[1]{#1}
\newcommand{\FloatTok}[1]{\textcolor[rgb]{0.25,0.63,0.44}{#1}}
\newcommand{\FunctionTok}[1]{\textcolor[rgb]{0.02,0.16,0.49}{#1}}
\newcommand{\ImportTok}[1]{\textcolor[rgb]{0.00,0.50,0.00}{\textbf{#1}}}
\newcommand{\InformationTok}[1]{\textcolor[rgb]{0.38,0.63,0.69}{\textbf{\textit{#1}}}}
\newcommand{\KeywordTok}[1]{\textcolor[rgb]{0.00,0.44,0.13}{\textbf{#1}}}
\newcommand{\NormalTok}[1]{#1}
\newcommand{\OperatorTok}[1]{\textcolor[rgb]{0.40,0.40,0.40}{#1}}
\newcommand{\OtherTok}[1]{\textcolor[rgb]{0.00,0.44,0.13}{#1}}
\newcommand{\PreprocessorTok}[1]{\textcolor[rgb]{0.74,0.48,0.00}{#1}}
\newcommand{\RegionMarkerTok}[1]{#1}
\newcommand{\SpecialCharTok}[1]{\textcolor[rgb]{0.25,0.44,0.63}{#1}}
\newcommand{\SpecialStringTok}[1]{\textcolor[rgb]{0.73,0.40,0.53}{#1}}
\newcommand{\StringTok}[1]{\textcolor[rgb]{0.25,0.44,0.63}{#1}}
\newcommand{\VariableTok}[1]{\textcolor[rgb]{0.10,0.09,0.49}{#1}}
\newcommand{\VerbatimStringTok}[1]{\textcolor[rgb]{0.25,0.44,0.63}{#1}}
\newcommand{\WarningTok}[1]{\textcolor[rgb]{0.38,0.63,0.69}{\textbf{\textit{#1}}}}
\setlength{\emergencystretch}{3em} % prevent overfull lines
\providecommand{\tightlist}{%
  \setlength{\itemsep}{0pt}\setlength{\parskip}{0pt}}
\setcounter{secnumdepth}{-\maxdimen} % remove section numbering
\usepackage{float}
\floatplacement{figure}{H}
\usepackage{bookmark}
\IfFileExists{xurl.sty}{\usepackage{xurl}}{} % add URL line breaks if available
\urlstyle{same}
\hypersetup{
  pdftitle={Advanced Machine Learning - Final Project},
  pdfauthor={Andres F. Camacho},
  colorlinks=true,
  linkcolor={blue},
  filecolor={Maroon},
  citecolor={blue},
  urlcolor={blue},
  pdfcreator={LaTeX via pandoc}}

\title{Advanced Machine Learning - Final Project}
\author{Andres F. Camacho}
\date{March 10, 2026}

\begin{document}
\maketitle

\subsection{Introduction}\label{introduction}

Last Saturday, March 7th, I participated with my team in the Harris
Public Policy Challenge (HPIC) on Chicago stadium and urban policy. The
question I faced was simple but daunting: how do you evaluate a policy
proposal when it's impossible to get everyone in Chicago in the same
room? The city has 2.7 million residents---each with their own
perspective on stadiums, subsidies, displacement, and civic identity.
Ignoring those voices isn't an option, but convening them all isn't
realistic either.

That's where the idea came from. After taking the course \emph{AI Agents
for Social Science} and \emph{Advanced Machine Learning for Public
Polocy}, I proposed using AI agents to simulate a deliberative
committee: expert panelists and community stakeholders who could
evaluate proposals together, push back on each other, and produce a
structured assessment. The goal wasn't to replace human judgment but to
\emph{inform} it---to surface tensions and tradeoffs that a single
analyst might miss, and to give voice to perspectives that are often
absent from policy debates. I ran my team's proposal through the system;
it didn't score high on all components, but the deliberation surfaced
gaps and tensions that helped us improve it before submission.

This report describes what I built, why, and how it works. Here's the
project in practice: you can explore the \textbf{live app} and the
\textbf{code} yourself.

\textbf{Try it:}
\href{https://ai-augmetnted-deliberative-committee-hpic.streamlit.app}{Streamlit
Dashboard} ·
\href{https://github.com/anfelipecb/AI-Augmetnted-Deliberative-Committee}{GitHub
Repository}

\subsection{The Problem}\label{the-problem}

The Harris challenge asks: \emph{What should be the City of Chicago's
policy toward supporting professional sports teams and facilities?}
Stadium deals are messy. They involve fiscal trade-offs (who pays, who
benefits?), equity concerns (who gets displaced, who gets jobs?), and
political feasibility (will City Council pass it?). Evaluating proposals
well means hearing from fiscal experts, community organizers, urban
economists, and residents who actually live near these projects. But
convening that panel is costly and slow. In practice, a lot of voices
never get heard.

I wanted to see if AI agents could help. Not to make the decision, but
to \emph{inform} it---by simulating a diverse panel that deliberates,
disagrees, and produces a structured evaluation. The idea: rather than
ignoring 2.7 million Chicagoans, I'd build agents that represent
different perspectives and let them debate.

\subsection{Literature Motivation and Recent
Developments}\label{literature-motivation-and-recent-developments}

This project is grounded in recent advancements in multi-agent LLM
systems, synthetic sampling, and AI ethics. The design of the
deliberative committee draws heavily from the following literature:

\textbf{Societies of Thought and Deliberation.} The core mechanism of
the committee is motivated by Kim, Evans, et al.~(2026), who demonstrate
that reasoning models generate ``societies of thought''---simulating
multi-agent-like dialogue to diversify perspectives, debate among
cognitive personas, and reconcile conflicting views. By structuring the
jury to deliberate across three rounds, this project harnesses these
emergent dynamics to surface tradeoffs that a single evaluator might
miss. Furthermore, Lai et al.~(2025) find that ``biased'' or opinionated
AI agents improve human decision-making by acting as strong heuristics;
similarly, assigning jury members explicit expert lenses (e.g., fiscal
constraint, urban economics) sharpens the debate and provides clear,
adversarial evaluation.

\textbf{Digital Doubles and Synthetic Subjects.} To capture community
impact, the pipeline simulates stakeholders using real demographic data.
This approach is motivated by Broska et al.~(2025), who propose the
``Mixed Subjects Design,'' treating LLMs as ``Silicon Subjects.'' This
framework validates using digital doubles to model population responses
when recruiting specific demographics is too costly. I sampled from ACS
2022 5-year estimates (income, housing, commute, occupation tables) to
ground each community profile in real demographic data for their
neighborhood and sector. The eight community stakeholders---a Pilsen
renter, an Austin business owner, a Bronzeville homeowner, and
others---are digital doubles across four geographic zones (South Side,
West Side, North Side, Loop-adjacent). Same proposal, same personas
yields reproducible evaluations; different proposals or modes let me
explore how perspectives shift.

\textbf{Persona Control and Ethics.} The distinct personalities and
evaluation criteria of the committee members are theoretically supported
by Chen et al.~(2025), who show that ``persona vectors'' can reliably
control character traits and analytical focuses in LLMs. Finally,
deploying such a system for urban policy requires careful consideration
of AI agency. Consistent with Gabriel et al.~(2025), who call for a new
ethics for a world of AI agents, this project emphasizes transparency:
deliberation logs are fully traceable, and the final synthesis preserves
both consensus and dissenting views, ensuring human oversight over the
policy evaluation process.

\subsection{How It Works}\label{how-it-works}

I use a two-tier setup: Opus for document understanding (summarizes the
proposal), Haiku for deliberation (cheaper, faster). Three modes:
\emph{jury only (quick)} with four experts, \emph{jury only (full)} with
seven, or \emph{full} with community stakeholders first, then the jury.

Three rounds. Round 1: each jury member scores Impact, Fiscal
Responsibility, Sustainability (1--10) plus six criteria
(LOW/MEDIUM/HIGH). Round 2: they see each other's scores and respond in
character---agreement, pushback, emphasis. Round 3: final scores,
two-sentence verdicts, and a synthesis report.

\textbf{Evaluation criteria.} The three main scores (Impact, Fiscal
Responsibility, Sustainability) and the six sub-criteria come from the
HPIC challenge's stated evaluation framework. The challenge asks whether
proposals benefit a cross-section of Chicagoans (Impact), generate tax
revenues and justify subsidies (Fiscal Responsibility), and support
sustainable, adaptive design (Sustainability). I operationalized these
into six criteria that jurors rate as LOW, MEDIUM, or HIGH:
\textbf{Fiscal Impact} (revenue, debt, public cost), \textbf{Equity \&
Access} (community benefit, affordability), \textbf{Political
Feasibility} (legal, legislative, stakeholder support),
\textbf{Sustainability} (long-term adaptability, risk), \textbf{Team
Retention} (likelihood of keeping major league franchises), and
\textbf{Accountability} (oversight, enforceability). These reflect the
policy considerations the challenge highlighted---public funding,
displacement, community engagement, and long-term fiscal
responsibility---so the deliberation aligns with what the City of
Chicago would evaluate.

\subsection{What I Built}\label{what-i-built}

Full pipeline (summarizer, jury rounds, community phase, synthesis),
seven jury personas and eight community personas, a Streamlit app, CLI,
and Docker. I also added an evaluation suite that validates outputs
(score ranges, criteria, verdict consistency) and a demo notebook with
alignment tests. Unit tests, integration test, and evaluation tests
round it out.

\subsection{Code Repository and Key
Components}\label{code-repository-and-key-components}

The project is organized as follows. Personas live in
\texttt{agents/jury/} and \texttt{agents/community/} (markdown files).
The orchestration logic is in \texttt{src/simulate.py}; agent invocation
and prompt building in \texttt{src/agents.py}; persona loading in
\texttt{src/personas.py}; criteria definitions in
\texttt{src/criteria.py}; and output formatting in
\texttt{src/output.py}. The evaluation suite is in
\texttt{src/evaluation.py}. The Streamlit app is \texttt{app.py}.

\textbf{Main simulation flow.} The \texttt{run()} function in
\texttt{src/simulate.py} orchestrates the full pipeline: summarizer
(Opus), optional community phase, then jury rounds 1--3 and synthesis.
Each round writes outputs to the run directory.

\begin{Shaded}
\begin{Highlighting}[]
\CommentTok{\# src/simulate.py — main run() flow (simplified)}
\KeywordTok{def}\NormalTok{ run(proposal\_text: }\BuiltInTok{str}\NormalTok{, mode: }\BuiltInTok{str} \OperatorTok{=} \StringTok{"jury"}\NormalTok{, output\_dir: Path }\OperatorTok{|} \VariableTok{None} \OperatorTok{=} \VariableTok{None}\NormalTok{) }\OperatorTok{{-}\textgreater{}}\NormalTok{ Path:}
\NormalTok{    proposal\_summary }\OperatorTok{=}\NormalTok{ \_run\_summarizer(proposal\_text, out)}
\NormalTok{    jury\_personas }\OperatorTok{=}\NormalTok{ load\_jury\_personas(quick}\OperatorTok{=}\NormalTok{(mode }\OperatorTok{==} \StringTok{"jury\_quick"}\NormalTok{))}
\NormalTok{    community\_summary }\OperatorTok{=} \StringTok{""}
    \ControlFlowTok{if}\NormalTok{ mode }\OperatorTok{==} \StringTok{"full"}\NormalTok{:}
\NormalTok{        community\_personas }\OperatorTok{=}\NormalTok{ load\_community\_personas()}
\NormalTok{        community\_summary, \_ }\OperatorTok{=}\NormalTok{ \_run\_community\_phase(proposal\_summary, community\_personas)}
\NormalTok{    round1\_scores, log\_entries }\OperatorTok{=}\NormalTok{ \_run\_round1\_jury(}
\NormalTok{        proposal\_summary, community\_summary, jury\_personas}
\NormalTok{    )}
\NormalTok{    log\_entries }\OperatorTok{=}\NormalTok{ \_run\_round2\_deliberation(jury\_personas, round1\_scores, log\_entries)}
\NormalTok{    final\_scores, synthesis, log\_entries }\OperatorTok{=}\NormalTok{ \_run\_round3\_final(jury\_personas, log\_entries, round1\_scores)}
    \CommentTok{\# ... write report.md, scores.json, deliberation\_log.md}
\end{Highlighting}
\end{Shaded}

\textbf{How prompts and context get passed.} Each agent receives a
\emph{system prompt} (persona + fixed criteria) and a \emph{user
message} (the proposal + context for that round). The system prompt is
built by concatenating the persona markdown with the evaluation
criteria:

\begin{Shaded}
\begin{Highlighting}[]
\CommentTok{\# src/agents.py — system prompt from persona + criteria}
\KeywordTok{def}\NormalTok{ build\_jury\_system\_prompt(persona\_content: }\BuiltInTok{str}\NormalTok{) }\OperatorTok{{-}\textgreater{}} \BuiltInTok{str}\NormalTok{:}
    \ControlFlowTok{return} \SpecialStringTok{f"""You are an expert panelist evaluating a Chicago stadium/urban policy proposal. Stay in character.}

\SpecialCharTok{\{}\NormalTok{persona\_content}\SpecialCharTok{\}}

\SpecialStringTok{You are evaluating the proposal against Chicago\textquotesingle{}s stated criteria. }\SpecialCharTok{\{}\NormalTok{CRITERIA\_TEXT}\SpecialCharTok{\}}
\SpecialCharTok{\{}\NormalTok{CRITERIA\_6\_TEXT}\SpecialCharTok{\}}

\SpecialStringTok{Respond in character... When asked for scores, you must respond with only a single valid JSON object..."""}
\end{Highlighting}
\end{Shaded}

\texttt{CRITERIA\_TEXT} injects the three main scores (Impact, Fiscal
Responsibility, Sustainability) with 1--10 scale and guidance;
\texttt{CRITERIA\_6\_TEXT} injects the six sub-criteria
(LOW/MEDIUM/HIGH) with sub-definitions from \texttt{src/criteria.py}.

\textbf{Round 1 user prompt.} The user message includes the proposal
summary (or community summary in full mode) and instructions for JSON
output. Each jury persona gets the same user message but a different
system prompt (their persona):

\begin{Shaded}
\begin{Highlighting}[]
\CommentTok{\# src/simulate.py — Round 1 prompt construction}
\NormalTok{prompt }\OperatorTok{=} \SpecialStringTok{f"""Below is a detailed summary of a Chicago stadium/urban policy proposal to evaluate. }\SpecialCharTok{\{}\NormalTok{intro}\SpecialCharTok{\}}

\SpecialStringTok{Proposal summary:}
\SpecialStringTok{{-}{-}{-}}
\SpecialCharTok{\{}\NormalTok{proposal\_text[:}\DecValTok{120000}\NormalTok{]}\SpecialCharTok{\}}
\SpecialStringTok{{-}{-}{-}}

\SpecialCharTok{\{}\NormalTok{community\_block}\SpecialCharTok{\}}
\SpecialStringTok{{-}{-}{-}}

\SpecialStringTok{Score the proposal. You must respond with only a single JSON object...}
\SpecialStringTok{"""}
\CommentTok{\# For each persona:}
\NormalTok{sys\_prompt }\OperatorTok{=}\NormalTok{ build\_jury\_system\_prompt(persona[}\StringTok{"content"}\NormalTok{])}
\NormalTok{response }\OperatorTok{=}\NormalTok{ invoke\_agent(persona[}\StringTok{"id"}\NormalTok{], sys\_prompt, prompt, }\VariableTok{None}\NormalTok{)}
\end{Highlighting}
\end{Shaded}

\textbf{Agent invocation.} The \texttt{invoke\_agent} function sends the
system prompt and messages to the Anthropic API. For Round 2 and 3,
\texttt{conversation\_history} can include prior rounds so the agent
sees the deliberation so far:

\begin{Shaded}
\begin{Highlighting}[]
\CommentTok{\# src/agents.py — API call}
\KeywordTok{def}\NormalTok{ invoke\_agent(agent\_id, system\_prompt, user\_message, conversation\_history}\OperatorTok{=}\VariableTok{None}\NormalTok{, model}\OperatorTok{=}\VariableTok{None}\NormalTok{):}
\NormalTok{    messages }\OperatorTok{=}\NormalTok{ []}
    \ControlFlowTok{if}\NormalTok{ conversation\_history:}
        \ControlFlowTok{for}\NormalTok{ m }\KeywordTok{in}\NormalTok{ conversation\_history:}
\NormalTok{            messages.append(\{}\StringTok{"role"}\NormalTok{: m[}\StringTok{"role"}\NormalTok{], }\StringTok{"content"}\NormalTok{: m[}\StringTok{"content"}\NormalTok{]\})}
\NormalTok{    messages.append(\{}\StringTok{"role"}\NormalTok{: }\StringTok{"user"}\NormalTok{, }\StringTok{"content"}\NormalTok{: user\_message\})}
\NormalTok{    response }\OperatorTok{=}\NormalTok{ client.messages.create(}
\NormalTok{        model}\OperatorTok{=}\NormalTok{model }\KeywordTok{or}\NormalTok{ DELIBERATION\_MODEL,}
\NormalTok{        max\_tokens}\OperatorTok{=}\DecValTok{4096}\NormalTok{,}
\NormalTok{        system}\OperatorTok{=}\NormalTok{system\_prompt,}
\NormalTok{        messages}\OperatorTok{=}\NormalTok{messages,}
\NormalTok{    )}
    \ControlFlowTok{return}\NormalTok{ response.content[}\DecValTok{0}\NormalTok{].text}
\end{Highlighting}
\end{Shaded}

So the flow is: persona content becomes the system prompt; the proposal
summary plus round-specific context becomes the user message; prior
rounds (when applicable) are passed as conversation\_history. The API
returns the agent's text, which is parsed for JSON scores and criteria.

\subsection{Verification}\label{verification}

The evaluation suite
(\texttt{uv\ run\ python\ -m\ src.evaluation\ -\/-run-dir\ outputs/\textless{}run\_id\textgreater{}})
checks that scores are valid, criteria are LOW/MEDIUM/HIGH, and verdicts
align with scores. Similar in spirit to the portfolio homework's
backtesting---I validate that the pipeline produces sensible, consistent
outputs.

\textbf{HPIC Demo Notebook.} I included
\texttt{docs/notebooks/hpic\_demo.ipynb} to showcase the
concept---allowing readers and the professor to explore the project
without running the full app. The notebook walks through: (1)
introduction and main components (summarization, deliberation, personas,
outputs); (2) simulation flow and the four quick jury profiles; (3) a
sample run with embedded data or loaded outputs from prior runs; (4)
visualizations (bar chart of scores by agent, criteria heatmap); (5)
alignment tests (persona differentiation, lens alignment---e.g., does
Sarah Chen weight fiscal highly? does Marcus Thompson weight impact and
equity\_access highly?); and (6) integration with the evaluation suite.
It works with or without prior runs thanks to embedded sample data, so
the concept is reproducible even when \texttt{outputs/} is empty.

\subsection{Effort}\label{effort}

I have been working on this project for the last couple of
months---investigating the literature on AI agents and digital doubles,
reading papers from the AI Agents for Social Science course, and trying
different combinations of models, prompts, and persona designs. The work
spans the full pipeline (summarization with Opus, deliberation with
Haiku), seven jury personas and eight community stakeholders grounded in
ACS 2022 data, the Streamlit app and CLI, Docker deployment, the
evaluation suite, and the demo notebook. I estimate roughly 50--70 hours
total: research and design (\textasciitilde15 h), persona development
and ACS data integration (\textasciitilde12 h), pipeline implementation
and debugging (\textasciitilde15 h), UI and deployment (\textasciitilde8
h), evaluation suite and documentation (\textasciitilde10 h). The HPIC
challenge on March 7th was a milestone; the Advanced ML final project
additions (evaluation suite, demo notebook, report) came afterward.

\subsection{Limitations}\label{limitations}

No formal hallucination detection yet. Persona differentiation is
heuristic. The verdict consistency check uses a fixed set of phrases.
Future work: claim extraction and grounding against the proposal, better
metrics for persona differentiation.

\subsection{Conclusion}\label{conclusion}

The HPIC AI-Augmented Deliberative Committee is an application of
multi-agent systems to policy deliberation. Digital doubles + societies
of thought = a way to surface diverse perspectives when you can't get
2.7 million people in a room. The project is on GitHub with full docs, a
demo notebook, and an evaluation suite. Feel free to run it, fork it, or
just explore the live app.

\textbf{Links:}
\href{https://ai-augmetnted-deliberative-committee-hpic.streamlit.app}{Streamlit
Dashboard} ·
\href{https://github.com/anfelipecb/AI-Augmetnted-Deliberative-Committee}{GitHub}

\subsection{Annex: Example Community Stakeholder
Profile}\label{annex-example-community-stakeholder-profile}

Below is one of the eight community stakeholder profiles (digital
doubles) used in full mode. Each profile is sampled from ACS 2022 tables
(e.g., B19013 income, B25064 rent, B25077 home values, B08303 commute)
and stored as a markdown file; the system prompt is built from this
content plus fixed criteria. This example---Rosa Delgado, a Pilsen
renter---illustrates the structure: personal profile, location and
housing, economic profile, community context, and stakeholder
perspective, with methodology notes citing the source tables.

\begin{Shaded}
\begin{Highlighting}[]
\FunctionTok{\# Rosa Delgado}

\FunctionTok{\#\# Personal Profile}
\SpecialStringTok{{-} }\NormalTok{Age: 34}
\SpecialStringTok{{-} }\NormalTok{Household composition: Married, 3 children (ages 12, 9, 6)}
\SpecialStringTok{{-} }\NormalTok{Years in Chicago/neighborhood: 16 years (grew up in Little Village, moved to Pilsen at 18)}
\SpecialStringTok{{-} }\NormalTok{Primary language: Spanish (limited English proficiency)}

\FunctionTok{\#\# Location \& Housing}
\SpecialStringTok{{-} }\NormalTok{Neighborhood: East Pilsen, near 18th \& Morgan}
\SpecialStringTok{{-} }\NormalTok{Housing tenure: Renter}
\SpecialStringTok{{-} }\NormalTok{Monthly rent: $1,350 (up from $950 three years ago)}

\FunctionTok{\#\# Economic Profile}
\SpecialStringTok{{-} }\NormalTok{Occupation: Restaurant server at downtown hotel}
\SpecialStringTok{{-} }\NormalTok{Annual household income: $47,800 (combined with husband\textquotesingle{}s construction work)}
\SpecialStringTok{{-} }\NormalTok{Education: High school diploma}
\SpecialStringTok{{-} }\NormalTok{Commute: 35{-}40 minutes via Pink Line to Loop}

\FunctionTok{\#\# Community Context}
\SpecialStringTok{{-} }\NormalTok{Neighborhood ties: Shops at Supermercado Monterrey on 18th St, children attend Benito Juárez Community Academy}
\SpecialStringTok{{-} }\NormalTok{Community involvement: Attends community meetings through church (St. Adalbert)}

\FunctionTok{\#\# Stakeholder Perspective}
\SpecialStringTok{{-} }\NormalTok{Primary concerns: Rent increases and displacement; transit reliability; affordable family activities; access to information in Spanish}
\SpecialStringTok{{-} }\NormalTok{Priorities: Affordable housing protections; living{-}wage jobs accessible by transit; family{-}friendly public spaces}
\SpecialStringTok{{-} }\NormalTok{Perspective on development: Skeptical of development without renter protections; worries about gentrification; needs concrete affordability guarantees}

\NormalTok{*Methodology: Income and rent from ACS 2022 Tables B19013, B25064 (Community Area 31 {-} Lower West Side/Pilsen)*}
\end{Highlighting}
\end{Shaded}

\subsection{References}\label{references}

All references were checked and browsed to verify titles, authors, and
URLs before inclusion, to prevent hallucinations.

\begin{itemize}
\tightlist
\item
  Kim, J., Lai, S., Scherrer, N., Agüera y Arcas, B., \& Evans, J.
  (2026). Reasoning Models Generate Societies of Thought.
  arXiv:2601.10825. https://arxiv.org/html/2601.10825v1
\item
  Broska et al.~(2025). The Mixed Subjects Design: Treating Large
  Language Models as Potentially Informative Observations.
\item
  Chen et al.~(2025). Persona Vectors: Monitoring and Controlling
  Character Traits in Language Models.
\item
  Lai et al.~(2025). Biased AI improves human decision-making but
  reduces trust.
\item
  Gabriel et al.~(2025). We need a new ethics for a world of AI agents.
\item
  U.S. Census Bureau. American Community Survey 2022 5-year estimates.
  https://data.census.gov
\item
  City of Chicago Data Portal. https://data.cityofchicago.org
\item
  Anthropic API. https://docs.anthropic.com
\end{itemize}

\end{document}
